\section{Дневник выполнения работы и анализ недочётов}

\subsection{Дневник выполнения работы}
\begin{enumerate}
    \item \textbf{День 1: Базовая проверка памяти.} Первоначальная реализация AVL-дерева была протестирована с помощью утилиты \texttt{leaks}. Благодаря корректной реализации рекурсивного метода \texttt{destroy()} в деструкторе, утечек памяти обнаружено не было (\texttt{0 leaks for 0 total leaked bytes}).

    Посмотрим на результаты неоптимизированных бенчмарков:

\begin{lstlisting}[language=bash]
artems@MacBook-Air-artem-3 lab2 % leaks --atExit -- ./benchmark_gprof_non_optim tests/01
benchmark_gprof(40739) MallocStackLogging: could not tag MSL-related memory as no_footprint, so those pages will be included in process footprint - (null)
benchmark_gprof(40739) MallocStackLogging: recording malloc and VM allocation stacks using lite mode
Loading data into memory... Done. Loaded 643408 operations.

Running std::map... Done.
Running avl::AVL... Done.

=== BENCHMARK RESULTS ===
Total Operations : 643408
-------------------------
std::map (RB)    :    790 ms | Mismatches: 0
avl::AVL         :   1379 ms | Mismatches: 0
-------------------------
SUCCESS: Both trees produced 100% correct results.
Process 40739 is not debuggable. Due to security restrictions, leaks can only show or save contents of readonly memory of restricted processes.

Process:         benchmark_gprof [40739]
Path:            /Users/USER/Desktop/*/benchmark_gprof
Load Address:    0x102f58000
Identifier:      benchmark_gprof
Version:         0
Code Type:       ARM64
Platform:        macOS
Parent Process:  leaks [40738]

Date/Time:       2026-04-12 18:20:57.219 +0300
Launch Time:     2026-04-12 18:20:52.600 +0300
OS Version:      macOS 14.6.1 (23G93)
Report Version:  7
Analysis Tool:   /usr/bin/leaks

Physical footprint:         263.8M
Physical footprint (peak):  547.4M
Idle exit:                  untracked
----

leaks Report Version: 4.0, multi-line stacks
Process 40739: 186 nodes malloced for 15 KB
Process 40739: 0 leaks for 0 total leaked bytes.

\end{lstlisting}

    \item \textbf{День 2: Профилирование времени.} Запуск профилировщика (\texttt{Samply}) показал, что программа работает медленнее эталонного \texttt{std::map}. Анализ графа вызовов выявил, что около 29\% процессорного времени тратится на файловый ввод (\texttt{std::istream}) и около 13\% на саму вставку в дерево. Также замечено значительное время, затрачиваемое на копирование ключей (\texttt{std::string}).
    \item \textbf{День 3: Оптимизация (Избавление от рекурсии).} Было принято решение отказаться от рекурсивного метода \texttt{insert}, так как каждый рекурсивный вызов создает новый фрейм стека. Написана итеративная версия вставки с использованием локального массива \texttt{Node* path[128]} для хранения пути. Ключи и значения начали передаваться через \texttt{std::move}.
    \item \textbf{День 4: Финальные замеры.} Проведено сравнительное тестирование неоптимизированной и оптимизированной версий программы под контролем \texttt{leaks}. 
\begin{lstlisting}[language=bash]
artems@MacBook-Air-artem-3 lab2 % g++ -std=c++17 -pg -O2 -g -Wall -Wextra \   
    benchmark.cpp -o benchmark_gprof
artems@MacBook-Air-artem-3 lab2 % leaks --atExit -- ./benchmark_gprof_optim tests/01
benchmark_gprof(41393) MallocStackLogging: could not tag MSL-related memory as no_footprint, so those pages will be included in process footprint - (null)
benchmark_gprof(41393) MallocStackLogging: recording malloc and VM allocation stacks using lite mode
Loading data into memory... Done. Loaded 643408 operations.

Running std::map... Done.
Running avl::AVL... Done.

=== BENCHMARK RESULTS ===
Total Operations : 643408
-------------------------
std::map (RB)    :    820 ms | Mismatches: 0
avl::AVL         :    782 ms | Mismatches: 0
-------------------------
SUCCESS: Both trees produced 100% correct results.
Process 41393 is not debuggable. Due to security restrictions, leaks can only show or save contents of readonly memory of restricted processes.

Process:         benchmark_gprof [41393]
Path:            /Users/USER/Desktop/*/benchmark_gprof
Load Address:    0x104df0000
Identifier:      benchmark_gprof
Version:         0
Code Type:       ARM64
Platform:        macOS
Parent Process:  leaks [41392]

Date/Time:       2026-04-12 18:24:04.060 +0300
Launch Time:     2026-04-12 18:24:00.439 +0300
OS Version:      macOS 14.6.1 (23G93)
Report Version:  7
Analysis Tool:   /usr/bin/leaks

Physical footprint:         9.8M
Physical footprint (peak):  296.0M
Idle exit:                  untracked
----

leaks Report Version: 4.0, multi-line stacks
Process 41393: 186 nodes malloced for 15 KB
Process 41393: 0 leaks for 0 total leaked bytes.
\end{lstlisting}

\end{enumerate}

\subsection{Выводы о найденных недочётах}
В ходе исследования были выявлены следующие архитектурные недочёты первой версии:
\begin{itemize}
    \item \textbf{Передача параметров по значению в рекурсии:} При передаче строк по значению (\texttt{std::string}) на каждом шаге рекурсии происходило выделение памяти под копию строки. Использование \texttt{const Key\&} или \texttt{std::move} критически важно для производительности.
    \item \textbf{Накладные расходы стека вызовов:} Рекурсивная вставка порождала глубокий стек. Это не только замедляло исполнение из-за сохранения регистров процессора, но и усложняло кэширование.
\end{itemize}

\subsection{Общие выводы об опыте}
Лабораторная работа наглядно продемонстрировала важность использования профилировщиков.
\newpage